\begin{afterword}
\summaryhead

The post draws on Robyn Dawes's \textsc{House of Cards}, which argues that research does not support psychotherapists' claims to superior insight and skill. The author gives examples. The Rorschach ink-blot test has failed hundreds of experiments and is still in use. Patients who see therapists get better faster than those who do nothing, but the schools of therapy do not differ in results, experience brings no gain, and talking to a college professor from another field works as well as seeing a therapist.

Even so, therapists were licensed, heard in court, trained in accredited schools and paid by insurers, all without experimental evidence. Schools multiplied, perhaps because no experiment could show one better than another. As far as the author can discern, prestige came from founding a school on charisma and good stories, as Freud had done. Happiness research, by contrast, began with measures that agree with each other. The lesson is that any organized practice needs measurement and controlled testing of its techniques. A later addition grants that matters may have improved and recalls evidence in favor of cognitive-behavioral therapy.

\responsehead

The post's main picture of psychotherapy research, as it stood in 2009, has support. Major reviews found little difference in outcome between schools of therapy, and paraprofessionals did about as well as professionals. The post's added line recalls, correctly, the evidence for cognitive-behavioral therapy that qualified the finding on schools. The contrast with happiness research also holds: its measures do agree with each other.

Beyond that picture, several claims are stated more strongly than the evidence allows. The post says the answer on the Rorschach is simply ``No.'' The most critical review before the post found support for ``a small number of indexes derived from the Rorschach and TAT'', a related test, and none for most of them. The college professors who did as well as therapists were not ``randomly selected''; in the one study the claim matches, they were ``chosen for their ability to form understanding relationships'', and the patients were male college students. From that study the post concludes, with ``apparently'', that talking to ``anyone'' helps. And the field is said to have been licensed and paid ``In the entire absence of the slightest experimental evidence for their effectiveness'', two paragraphs after the post grants that patients in therapy get better faster than those who do nothing. The post gives no dates that would let the sentence refer to an earlier time.

The explanation for the proliferation of schools, that prestige came from founding one, is offered with ``So far as I can discern'' and ``probably''. It is a hypothesis, and the post gives no evidence for it.

In short: a picture of psychotherapy research that the reviews of the time support, with some claims stated more strongly than the studies allow (a flat ``No'' for the Rorschach, ``randomly selected'' professors, ``the entire absence'' of evidence).
\end{afterword}
